\cotefile{97}
\scanpage{2}
% letterhead
\noindent\begin{tabular}[t]{c}
COLLÈGE\\
DE\\
FRANCE\\
\rule{1.8em}{0.5pt}\\
CHAIRE\\
D’ALGÈBRE ET GÉOMÉTRIE
\end{tabular}\hfill
\begin{tabular}[t]{l}
\textit{Paris, le} 4 Nov. 66\\
\end{tabular}

Cher Grothendieck,\\

Je n’ai pas regardé ton pb. Mais je suis tombé sur une ``Notice'' que je t’envoie ci-joint, démontrant la même chose que toi.

La Notice renvoie à un papier de Milnor, que j’ai, où celui-ci regarde justement \unsure{le fibré|les fibrés} à groupe $\mathrm{SL}_{2}$ qui sont plats, et de base une surface. C’est \unsure{très} joli.

Bien à toi\\
J-P Se\uline{rr}e.\\
% archive sheet number: 1
\clearpage
